\documentclass[12pt]{article}
\usepackage[utf8]{inputenc}
\usepackage[T1]{fontenc}
\usepackage[margin=1in]{geometry}
\usepackage{setspace}
\usepackage{xcolor}
\usepackage[colorlinks=true, urlcolor=blue!60!black, linkcolor=blue!60!black]{hyperref}
\setlength{\parindent}{0em}
\setlength{\parskip}{0.45em}
\renewcommand{\baselinestretch}{1.0}
\begin{document}
\begin{center}
{\Large \textit{Ashcroft v.\ Free Speech Coalition}, 535 U.S.\ 234 (2002)\par}
\vspace{0.3em}
{\normalsize Case Brief --- One-Page Edition\par}
\end{center}

\textbf{Memory Jogger.} Virtual child pornography: non-obscene speech made without harming real children is protected (\href{https://supreme.justia.com/cases/federal/us/535/234/}{535 U.S.\ 234}).

\textbf{Facts.} The Child Pornography Prevention Act of 1996 banned images that ``appear to be'' of minors in sexually explicit conduct, reaching computer-generated images and youthful-looking adults, and depictions marketed so as to ``convey the impression'' of minors. An adult-entertainment trade group and artists sued.

\textbf{Procedural History.} The district court upheld the Act; the Ninth Circuit reversed; certiorari was granted.

\textbf{Votes.} Affirmed, 6--3 on virtual images and 7--2 on youthful adults and pandering (Kennedy for five; Thomas concurring in the judgment; O'Connor in part; Rehnquist and Scalia dissenting).

\textbf{Holding.} Both provisions are facially overbroad. Speech that is neither obscene under \textit{Miller} nor produced by exploiting real children under \textit{Ferber} keeps its protection, and Congress may not ban it because it might indirectly encourage abuse or complicate prosecutions.

\textbf{Reasoning.} The Act ignored \textit{Miller}'s prongs, threatening works like \textit{Romeo and Juliet}, and \textit{Ferber} rested on speech ``intrinsically related'' to real abuse. The mere tendency of speech to encourage unlawful thoughts cannot justify suppression, and banning protected speech because it resembles unprotected speech ``turns the First Amendment upside down.'' The pandering clause tainted innocent possessors. Thomas left room for a narrow future law; O'Connor and Rehnquist would have upheld a ban on images ``virtually indistinguishable from'' real children.

\textbf{Examples.} \emph{Same-Side:} a ban on non-obscene animated teenagers falls. \emph{Opposite-Side:} prosecuting obscene virtual images under \textit{Miller}. \emph{Fence-Sitter:} generative-AI images from models trained on real abuse material.

\textbf{Critique.} Child-protection scholars say the production-harm test understates real dangers, especially with AI; doctrinal critics fault the use of overbreadth over narrowing construction.

\textbf{Key Quotation.} ``In contrast to the speech in \textit{Ferber}, speech that itself is the record of sexual abuse, the CPPA prohibits speech that records no crime and creates no victims by its production. Virtual child pornography is not `intrinsically related' to the sexual abuse of children, as were the materials in \textit{Ferber}.'' --- Justice Kennedy, 535 U.S.\ at 250.

\end{document}
